% Laget automatisk av elevark/lag_elevark.py fra opplegg/kvanteobjekter.html. Endringer her blir overskrevet.
\documentclass[11pt]{article}
\usepackage{elevark}
\begin{document}
\elevarkhode{Fysikk 2 · Kvantefysikk}{Bølger og partikler på samme tid – arbeidsark}{https://fysikklab.no/sim/dobbeltspalte.html}{Dobbeltspalte}

\section*{Oppgaver}

Oppgave 1–4 bruker Dobbeltspalte med én partikkel om gangen\footnote{\url{fysikklab.no/sim/dobbeltspalte.html}}. Bruk $h = 6{,}63 \cdot 10^{-34}$ Js, $m_{\mathrm{e}} = 9{,}11 \cdot 10^{-31}\,\mathrm{kg}$ og $c = 3{,}00 \cdot 10^{8}\,\mathrm{m/s}$.

\begin{oppgaveliste}
\oppgave \fase{gjett}{Gjett}Klinkekuler sendes mot to spalter. Hvordan ser mønsteret ut? Prøv, og bytt så til fotoner med 10 per sekund.
\svar
\oppgave \fase{forklar}{Forklar}Hvert foton treffer ett punkt. Hvordan kan det likevel bli et stripemønster etter mange fotoner?
\svar
\oppgave \fase{test}{Undersøk}Velg «To spalter, mål hvilken spalte». Hva skjer med mønsteret? Sammenlign med én spalte.
\svar
\oppgave \fase{test}{Undersøk}Velg elektroner. Hva skjer med stripene når du øker akselerasjonsspenningen? Forklar med de Broglie-bølgelengden.
\svar
\oppgave \fase{test}{Regn}Et røntgenfoton med bølgelengde $1{,}43 \cdot 10^{-12}\,\mathrm{m}$ spres 90° i et comptonstøt. Hvor stor blir bølgelengden etterpå?
\svar
\oppgave \fase{test}{Regn}Finn de Broglie-bølgelengden til et elektron som er akselerert gjennom 100~V, og til en klinkekule på 21~g med farten 5,0~m/s.
\svar
\oppgave \fase{test}{Regn}Et elektron er fanget i et hydrogenatom med diameter $1{,}0 \cdot 10^{-10}\,\mathrm{m}$. Finn den minste uskarpheten i bevegelsesmengden, og den tilsvarende kinetiske energien i eV.
\svar
\oppgave \fase{utfordring}{Utfordring}Et elektronmikroskop skal se et ribosom på 20~nm. Bølgelengden må være mindre enn det. Hvor stor bevegelsesmengde og kinetisk energi må elektronene minst ha? Hvorfor kan ikke et lysmikroskop se ribosomet?
\svar
\end{oppgaveliste}

\end{document}
